\section{Find a Peak Element in a 2D Matrix}
\label{sec:bs-peak-2d}

\problemstatement{We are given an $n \times m$ matrix of distinct
integers, where no two horizontally or vertically adjacent cells are
equal. A cell is a \textbf{peak} if it is strictly greater than all of its
(up to four) horizontal and vertical neighbors; cells outside the matrix
are treated as $-\infty$. Find the position of any one peak.}

\begin{patternbox}
This is the 2D generalization of Section~\ref{sec:bs-peak-1d}, and the
keyword pattern is identical: no sortedness, just a local-comparison rule
plus an implicit $-\infty$ border. The new idea needed is how to binary
search over \emph{one} dimension (say, columns) while using an $O(n)$
linear scan over the \emph{other} dimension (rows) as part of each
predicate evaluation -- a ``binary search over one axis, brute-force scan
over the other'' technique that is common whenever a 2D problem reduces
to a 1D binary search plus a per-step linear helper.
\end{patternbox}

\subsection*{Understanding the problem carefully}

Fix a candidate middle column $\textit{midCol}$. Scan down that column to
find the row $\textit{maxRow}$ containing the largest value in that
column -- call it $v = \textit{matrix}[\textit{maxRow}][\textit{midCol}]$.
Compare $v$ to its left neighbor $\textit{matrix}[\textit{maxRow}]
[\textit{midCol}-1]$ and right neighbor
$\textit{matrix}[\textit{maxRow}][\textit{midCol}+1]$ (treating
out-of-range neighbors as $-\infty$).

\begin{itemize}
  \item If $v$ is greater than both neighbors: $v$ is greater than
        everything above and below it in its own column (by choice of
        $\textit{maxRow}$ as the column's maximum) \emph{and} greater than
        its left and right neighbors -- so $v$ is a peak. Done.
  \item If the left neighbor is greater than $v$: since that left
        neighbor is, in particular, greater than the column maximum at
        $\textit{midCol}$, the column $\textit{midCol}-1$ (or further
        left) must contain an even larger local structure -- a peak is
        guaranteed to exist somewhere in the columns
        $[\textit{lo}, \textit{midCol}-1]$, by the same $-\infty$-border
        existence argument as Section~\ref{sec:bs-peak-1d}, applied along
        the column axis.
  \item Symmetrically, if the right neighbor is greater, a peak is
        guaranteed in columns $[\textit{midCol}+1, \textit{hi}]$.
\end{itemize}

\begin{ideabox}[Candidate Space and Predicate]
The candidate space is columns $0, \dots, m-1$. At each candidate column,
an $O(n)$ scan finds that column's maximum value's row, and a constant-time
comparison to its horizontal neighbors decides whether to stop (peak
found) or which half of the remaining columns to search next -- the same
``climbing toward a guaranteed peak'' argument as
Section~\ref{sec:bs-peak-1d}, just with an entire column's maximum
standing in for a single array element.
\end{ideabox}

\subsection*{Why restricting attention to each column's maximum is valid}

The reason we only need to examine the \emph{maximum} of each candidate
column, rather than every cell in it, is that if a peak exists in this
column, it is impossible for it to be at a position other than the
column's maximum: a peak must be greater than its upper and lower
neighbors within the same column, which already forces it to be at least
as large as both of those, and chaining this reasoning shows the only cell
in a column that could possibly satisfy ``greater than both vertical
neighbors'' even before checking horizontal neighbors is necessarily a
local maximum along that column. Examining the column's global maximum is
therefore sufficient to either confirm a peak or correctly decide which
direction to continue the binary search.

\subsection*{Algorithm}

\begin{enumerate}
  \item Initialize $\textit{lo} \gets 0$, $\textit{hi} \gets m-1$.
  \item While $\textit{lo} \le \textit{hi}$:
  \begin{enumerate}
    \item $\textit{midCol} \gets \textit{lo}+(\textit{hi}-\textit{lo})/2$.
    \item Scan column $\textit{midCol}$ to find $\textit{maxRow}$, the row
          of its maximum value $v$.
    \item Compare $v$ to its left and right neighbors.
    \item If $v$ exceeds both: return $(\textit{maxRow},
          \textit{midCol})$.
    \item If the left neighbor exceeds $v$: $\textit{hi} \gets
          \textit{midCol}-1$.
    \item Else (right neighbor exceeds $v$): $\textit{lo} \gets
          \textit{midCol}+1$.
  \end{enumerate}
\end{enumerate}

\begin{algorithm}[H]
\caption{Find Peak Element in a 2D Matrix}
\begin{algorithmic}[1]
\Procedure{FindPeak2D}{\textit{matrix}}
  \State $\textit{lo} \gets 0$, $\textit{hi} \gets m-1$
  \While{$\textit{lo} \le \textit{hi}$}
    \State $\textit{midCol} \gets \textit{lo}+(\textit{hi}-\textit{lo})/2$
    \State $\textit{maxRow} \gets \arg\max_i \textit{matrix}[i][\textit{midCol}]$
    \State $v \gets \textit{matrix}[\textit{maxRow}][\textit{midCol}]$
    \State $\textit{left} \gets (\textit{midCol}=0)\ ? -\infty : \textit{matrix}[\textit{maxRow}][\textit{midCol}-1]$
    \State $\textit{right} \gets (\textit{midCol}=m-1)\ ? -\infty : \textit{matrix}[\textit{maxRow}][\textit{midCol}+1]$
    \If{$v > \textit{left}$ \textbf{and} $v > \textit{right}$}
      \State \Return $(\textit{maxRow}, \textit{midCol})$
    \ElsIf{$\textit{left} > v$}
      \State $\textit{hi} \gets \textit{midCol} - 1$
    \Else
      \State $\textit{lo} \gets \textit{midCol} + 1$
    \EndIf
  \EndWhile
\EndProcedure
\end{algorithmic}
\end{algorithm}

\subsection*{Dry run}

\dryrun{Take $\textit{matrix} = \begin{bmatrix} 10&8&10&10 \\ 14&13&12&11
\\ 15&9&11&21 \\ 16&17&19&20 \end{bmatrix}$ ($n=m=4$).}

\begin{itemize}
  \item $\textit{lo}=0,\textit{hi}=3$: $\textit{midCol}=1$. Column $1$:
        $[8,13,9,17]$; max is $17$ at row $3$. Left neighbor
        $\textit{matrix}[3][0]=16$, right neighbor
        $\textit{matrix}[3][2]=19$. Is $17 > 16$ and $17 > 19$? No
        ($17 < 19$). Right neighbor is greater: $\textit{lo}=2$.
  \item $\textit{lo}=2,\textit{hi}=3$: $\textit{midCol}=2$. Column $2$:
        $[10,12,11,19]$; max is $19$ at row $3$. Left neighbor
        $\textit{matrix}[3][1]=17$, right neighbor
        $\textit{matrix}[3][3]=20$. Is $19>17$ and $19>20$? No
        ($19<20$). Right neighbor greater: $\textit{lo}=3$.
  \item $\textit{lo}=3,\textit{hi}=3$: $\textit{midCol}=3$. Column $3$:
        $[10,11,21,20]$; max is $21$ at row $2$. Left neighbor
        $\textit{matrix}[2][2]=11$. No right neighbor (out of range,
        $-\infty$). Is $21>11$ and $21>-\infty$? Yes. Peak found at
        $(2,3)$.
\end{itemize}

\subsection*{C++ implementation}

\begin{lstlisting}
#include <bits/stdc++.h>
using namespace std;

pair<int,int> findPeak2D(vector<vector<int>>& matrix) {
    int n = matrix.size(), m = matrix[0].size();
    int lo = 0, hi = m - 1;

    while (lo <= hi) {
        int midCol = lo + (hi - lo) / 2;
        int maxRow = 0;
        for (int i = 1; i < n; i++) {
            if (matrix[i][midCol] > matrix[maxRow][midCol]) maxRow = i;
        }

        int v = matrix[maxRow][midCol];
        int left  = (midCol == 0)     ? INT_MIN : matrix[maxRow][midCol - 1];
        int right = (midCol == m - 1) ? INT_MIN : matrix[maxRow][midCol + 1];

        if (v > left && v > right) {
            return {maxRow, midCol};
        } else if (left > v) {
            hi = midCol - 1;
        } else {
            lo = midCol + 1;
        }
    }
    return {-1, -1}; // unreachable for a valid matrix
}

int main() {
    vector<vector<int>> matrix = {
        {10, 8, 10, 10},
        {14, 13, 12, 11},
        {15, 9, 11, 21},
        {16, 17, 19, 20}
    };
    auto [r, c] = findPeak2D(matrix);
    cout << "Peak at (" << r << ", " << c << ")" << endl;
    return 0;
}
\end{lstlisting}

\begin{complexitybox}
\textbf{Time Complexity.} The column range halves each iteration
($O(\log m)$ iterations), and each iteration scans a full column in
$O(n)$, giving an overall time complexity of
\[
  O(n \log m).
\]
\textbf{Space Complexity.} $O(1)$ auxiliary space.
\end{complexitybox}

\begin{takeawaybox}
When a 2D problem's local-comparison structure generalizes a 1D pattern
you already know (here, peak-finding from
Section~\ref{sec:bs-peak-1d}), look for a way to binary search one axis
while reducing each candidate to a single representative value via a
linear scan along the other axis -- the representative value's own local
comparisons then drive the binary search exactly as in the 1D version.
\end{takeawaybox}
